\chapter{Determinants: Volume and Sign}

In Chapter~20, Gram--Schmidt straightened columns into perpendicular altitudes whose product
measured the volume of a box. You already know that a matrix moves all of space at once --
stretching, rotating, shearing, and sometimes flattening. We now ask: \emph{by how much does the machine inflate or
deflate space?}

Feed the machine a little square of area $1$. It hands you back a parallelogram.
Measure that parallelogram's area. That single number -- the area-scaling factor
of the transformation -- is almost the whole story of the determinant. We say
``almost'' because there is one more subtlety: the determinant also
records whether the transformation reflects space, as in a mirror. So the
determinant carries \emph{two} pieces of geometric information at once: \textbf{how much
volume changed}, and \textbf{whether orientation was preserved or reversed}.

Every algebraic rule in this chapter -- the $2\times 2$ formula $ad-bc$, the
$3\times 3$ expansion, the condition $\det = 0$ for collapse, and the product law
$\det(\mat{A}\mat{B}) = \det\mat{A}\,\det\mat{B}$ -- follows directly from tracking
how a matrix scales volume and whether it flips orientation.

\section{The Unit Box and Its Image}
\label{ch21:unit-box}

Start in the plane with the cleanest possible region: the unit square. Its two
sides are the standard basis vectors,
\[
  \vec{e}_1 = \colvec{1\\0}, \qquad \vec{e}_2 = \colvec{0\\1},
\]
and it has area exactly $1$. Now act with a $2\times 2$ matrix $\mat{A}$. The
square has corners $\vec{0}$, $\vec{e}_1$, $\vec{e}_1+\vec{e}_2$, $\vec{e}_2$, and
$\mat{A}$ maps each corner to a new point. Because $\mat{A}\vec{e}_1$ is the
first column of $\mat{A}$ and $\mat{A}\vec{e}_2$ is the second column, the image
of the unit square is the parallelogram spanned by the two \textbf{columns} of
$\mat{A}$.

\begin{figure}[H]
\centering
\begin{tikzpicture}[scale=1.0,>=Stealth]
  % left: unit square
  \begin{scope}[xshift=0cm]
    \draw[->,gray!60] (-0.4,0) -- (2.6,0) node[right]{$x$};
    \draw[->,gray!60] (0,-0.4) -- (0,2.6) node[above]{$y$};
    \fill[vecblue!18] (0,0) -- (1,0) -- (1,1) -- (0,1) -- cycle;
    \draw[vecblue,very thick,->] (0,0) -- (1,0) node[midway,below=2pt]{$\vec{e}_1$};
    \draw[vecblue,very thick,->] (0,0) -- (0,1) node[midway,left=2pt]{$\vec{e}_2$};
    \node at (0.5,0.5) {\small area $1$};
    \node[below] at (0.5,-0.7) {unit square};
  \end{scope}
  % arrow
  \node at (3.5,1.0) {\large $\xrightarrow{\;\mat{A}\;}$};
  % right: parallelogram
  \begin{scope}[xshift=5cm]
    \draw[->,gray!60] (-0.4,0) -- (3.4,0) node[right]{$x$};
    \draw[->,gray!60] (0,-0.4) -- (0,2.6) node[above]{$y$};
    \fill[targetgreen!18] (0,0) -- (2,0.5) -- (2.5,1.7) -- (0.5,1.2) -- cycle;
    \draw[targetgreen,very thick,->] (0,0) -- (2,0.5) node[right=2pt]{$\mat{A}\vec{e}_1$};
    \draw[targetgreen,very thick,->] (0,0) -- (0.5,1.2) node[midway,left=2pt]{$\mat{A}\vec{e}_2$};
    \node at (1.25,0.85) {\small area $=\abs{\det\mat{A}}$};
    \node[below] at (1.25,-0.7) {image parallelogram};
  \end{scope}
\end{tikzpicture}
\caption{The determinant is the area-scaling factor. The machine $\mat{A}$ sends
the unit square (area $1$) to the parallelogram spanned by its columns
$\mat{A}\vec{e}_1$ and $\mat{A}\vec{e}_2$. The area of that parallelogram is
$\abs{\det\mat{A}}$.}
\label{ch21:fig-unitbox}
\end{figure}

\begin{intuitionbox}
The determinant of $\mat{A}$ is the \textbf{area of the parallelogram you get by
sending the unit square through $\mat{A}$} -- with a sign attached. The columns of
$\mat{A}$ are the two sides of that parallelogram. A determinant of $5$ means the
machine multiplies every area by $5$; a determinant of $\tfrac{1}{2}$ means it
halves every area; a determinant of $0$ means it flattens the square to a sliver
of zero area.
\end{intuitionbox}

Why does this scaling apply to every region, and not just the unit square? Because the machine is
\emph{linear}, and linearity is exactly the property that the grid stays a grid.
If the unit square's area is multiplied by some factor $k$, then so is the area of
the square next to it, and the one above it, and a $1000\times 1000$ block of
them tiled together. Any region you like can be approximated by tiny squares of
the grid; each one gets its area scaled by $k$; so the whole region's area gets
scaled by $k$. The determinant is one number that controls \emph{all} areas at
once. That is what makes it worth a name.

\begin{focusbox}
Hold onto three sentences. (1) $\det\mat{A}$ is the signed area-scaling factor of
the machine $\mat{A}$. (2) The columns of $\mat{A}$ are the edges of the image of
the unit box. (3) It scales \emph{every} region by the same factor, because
linear maps treat all of space uniformly.
\end{focusbox}

\section{The $2\times 2$ Determinant: Deriving \texorpdfstring{$ad-bc$}{ad-bc}}
\label{ch21:2x2}

We met the formula $ad-bc$ as an algebraic condition for $2\times 2$ invertibility
back in Chapter~6 and Chapter~8. Now let us see why it had to be that exact
combination of letters by computing the area of the column parallelogram directly.
Take a general $2\times 2$ matrix and read off its columns:
\[
  \mat{A} = \begin{bmatrix} a & b \\ c & d \end{bmatrix},
  \qquad
  \vec{u} = \mat{A}\vec{e}_1 = \colvec{a\\c},
  \qquad
  \vec{v} = \mat{A}\vec{e}_2 = \colvec{b\\d}.
\]
We want the area of the parallelogram with sides $\vec{u}$ and $\vec{v}$. The
cleanest geometric derivation is to enclose the tilted parallelogram inside a
bounding box of width $a+b$ and height $c+d$, then subtract the pieces that lie
outside the tilt (Figure~\ref{ch21:fig-shear-area}).

\begin{figure}[H]
\centering
\begin{tikzpicture}[scale=1.35,>=Stealth]
  % coordinates: a=3.2, c=1.0, b=1.2, d=2.4 -> (a+b)=4.4, (c+d)=3.4
  \coordinate (O) at (0,0);
  \coordinate (A) at (3.2,0);
  \coordinate (AB) at (4.4,0);
  \coordinate (U) at (3.2,1.0);
  \coordinate (UV) at (4.4,3.4);
  \coordinate (V) at (1.2,2.4);
  \coordinate (Btop) at (1.2,3.4);
  \coordinate (Cleft) at (0,2.4);
  \coordinate (Ctop) at (0,3.4);
  \coordinate (Uright) at (4.4,1.0);

  % 1. Surrounding regions (soft fills with clear labels)
  % bottom right triangle: base a, height c
  \fill[vecblue!15] (O) -- (A) -- (U) -- cycle;
  \draw[gray!50] (O) -- (A) -- (U) -- cycle;
  \node[font=\footnotesize,vecblue!80!black] at (2.2,0.35) {$\frac{1}{2}ac$};

  % bottom right rectangle: width b, height c
  \fill[gray!18] (A) -- (AB) -- (Uright) -- (U) -- cycle;
  \draw[gray!50] (A) -- (AB) -- (Uright) -- (U) -- cycle;
  \node[font=\footnotesize,gray!80!black] at (3.8,0.5) {$bc$};

  % top right triangle: base b, height d
  \fill[quotientorange!20] (U) -- (Uright) -- (UV) -- cycle;
  \draw[gray!50] (U) -- (Uright) -- (UV) -- cycle;
  \node[font=\footnotesize,quotientorange!90!black] at (4.0,1.8) {$\frac{1}{2}bd$};

  % top left triangle: base a, height c
  \fill[vecblue!15] (V) -- (Btop) -- (UV) -- cycle;
  \draw[gray!50] (V) -- (Btop) -- (UV) -- cycle;
  \node[font=\footnotesize,vecblue!80!black] at (2.6,3.05) {$\frac{1}{2}ac$};

  % top left rectangle: width b, height c
  \fill[gray!18] (Cleft) -- (O|-Ctop) -- (Btop) -- (V) -- cycle;
  \draw[gray!50] (Cleft) -- (O|-Ctop) -- (Btop) -- (V) -- cycle;
  \node[font=\footnotesize,gray!80!black] at (0.6,2.9) {$bc$};

  % bottom left triangle: base b, height d
  \fill[quotientorange!20] (O) -- (Cleft) -- (V) -- cycle;
  \draw[gray!50] (O) -- (Cleft) -- (V) -- cycle;
  \node[font=\footnotesize,quotientorange!90!black] at (0.4,1.55) {$\frac{1}{2}bd$};

  % 2. Central parallelogram
  \fill[targetgreen!22] (O) -- (U) -- (UV) -- (V) -- cycle;
  \draw[targetgreen,very thick] (O) -- (U) -- (UV) -- (V) -- cycle;
  \node[targetgreen!70!black,font=\small\bfseries] at (2.42,1.95) {Area $= ad - bc$};

  % 3. Enclosing bounding box
  \draw[gray!65,thick,dashed] (O) rectangle (UV);

  % 4. Vectors u and v
  \draw[vecblue,line width=1.3pt,->] (O) -- (U) node[pos=0.48,above=4pt,font=\footnotesize] {$\vec{u}=(a,c)$};
  \draw[quotientorange,line width=1.3pt,->] (O) -- (V) node[pos=0.44,right=4pt,font=\footnotesize] {$\vec{v}=(b,d)$};

  % 5. Dimension labels
  \draw[<->,gray!70] (0,-0.35) -- (3.2,-0.35) node[midway,below,font=\footnotesize] {$a$};
  \draw[<->,gray!70] (3.2,-0.35) -- (4.4,-0.35) node[midway,below,font=\footnotesize] {$b$};
  \draw[<->,gray!70] (4.75,0) -- (4.75,1.0) node[midway,right,font=\footnotesize] {$c$};
  \draw[<->,gray!70] (4.75,1.0) -- (4.75,3.4) node[midway,right,font=\footnotesize] {$d$};
\end{tikzpicture}
\caption{The geometric dissection derivation of $\det\mat{A}=ad-bc$. Enclose the
parallelogram spanned by $\vec{u}=(a,c)$ and $\vec{v}=(b,d)$ in a rectangle of size
$(a+b)\times(c+d)$. Subtracting the two blue triangles of area $\frac{1}{2}ac$, the
two orange triangles of area $\frac{1}{2}bd$, and the two gray corner rectangles of
area $bc$ leaves exactly $(ac+ad+bc+bd) - (ac+bd+2bc) = ad-bc$.}
\label{ch21:fig-shear-area}
\end{figure}

The total area of the bounding rectangle is:
\[
  \text{Total Box Area} = (a+b)(c+d) = ac + ad + bc + bd.
\]
Subtract the six pieces that fall outside the parallelogram:
\begin{itemize}
  \item Two right triangles with base $a$ and height $c$: combined area $2\times\big(\tfrac{1}{2}ac\big) = ac$.
  \item Two right triangles with base $b$ and height $d$: combined area $2\times\big(\tfrac{1}{2}bd\big) = bd$.
  \item Two corner rectangles of width $b$ and height $c$: combined area $2\times(bc) = 2bc$.
\end{itemize}
Subtracting these outer regions from the bounding box leaves the area of the parallelogram:
\[
  \text{Area} = (ac + ad + bc + bd) - (ac + bd + 2bc) = ad - bc.
\]
This is the $2\times 2$ formula, not as a symbol-pushing rule, but as the natural
remainder when the surrounding frame is removed.

We can also reach the exact same formula through base and height using the dot product:

\begin{intuitionbox}
The signed area of the parallelogram with edge vectors $\vec{u}=(u_1,u_2)$ and
$\vec{v}=(v_1,v_2)$ is $u_1 v_2 - u_2 v_1$. One quick way to remember why: area is
``base times height,'' and after you write base $=\norm{\vec u}$ and height
$=\norm{\vec v}\sin\theta$ (the part of $\vec v$ perpendicular to $\vec u$), the
product $\norm{\vec u}\,\norm{\vec v}\sin\theta$ expands exactly to
$u_1 v_2 - u_2 v_1$.
\end{intuitionbox}

Let us check that algebra directly. The area of a parallelogram is
base times height. Take the base to be $\vec u$, with length $\norm{\vec u}$. The
height is the component of $\vec v$ perpendicular to $\vec u$, namely
$\norm{\vec v}\sin\theta$ where $\theta$ is the angle from $\vec u$ to $\vec v$.
So
\[
  \text{area} = \norm{\vec u}\,\norm{\vec v}\,\sin\theta.
\]
Now recall that the dot product gives $\norm{\vec u}\,\norm{\vec v}\cos\theta =
u_1 v_1 + u_2 v_2$. There is a perpendicular twin: if you rotate $\vec u$
counterclockwise by $90^\circ$ to get $\vec u^{\perp} = (-u_2, u_1)$, then the
angle from $\vec u^\perp$ to $\vec v$ is $\theta - \tfrac{\pi}{2}$, so
\[
  \vec u^{\perp}\cdot \vec v
  = \norm{\vec u^\perp}\,\norm{\vec v}\,\cos\!\big(\theta-\tfrac{\pi}{2}\big)
  = \norm{\vec u}\,\norm{\vec v}\,\sin\theta = \text{signed area},
\]
using $\norm{\vec u^\perp} = \norm{\vec u}$ and
$\cos(\theta-\tfrac{\pi}{2}) = \sin\theta$. Computing the same dot product
component by component, $\vec u^{\perp}\cdot \vec v = -u_2 v_1 + u_1 v_2 =
u_1 v_2 - u_2 v_1$. When $\vec v$ lies counterclockwise from $\vec u$ this is
positive; when it lies clockwise it is negative -- the sign is the orientation we
will study in Section~\ref{ch21:sign}. Plugging in our columns
$\vec u = (a,c)$ and $\vec v = (b,d)$ (Figure~\ref{ch21:fig-base-height}):
\[
  \boxed{\ \det\!\begin{bmatrix} a & b \\ c & d \end{bmatrix} = ad - bc.\ }
\]

\begin{figure}[H]
\centering
\begin{tikzpicture}[scale=1.2,>=Stealth]
  \coordinate (O) at (0,0);
  \coordinate (U) at (3.2,0.6);
  \coordinate (V) at (1.0,2.2);
  \coordinate (UV) at ($(U)+(V)$);
  \coordinate (P) at (1.365, 0.256);

  \fill[targetgreen!18] (O) -- (U) -- (UV) -- (V) -- cycle;
  \draw[targetgreen,thick] (O) -- (U) -- (UV) -- (V) -- cycle;

  \draw[dashed,gray!75,thick] (V) -- (P) node[midway,left=2pt,font=\footnotesize,gray!85!black] {$h$};

  \draw[vecblue,very thick,->] (O) -- (U) node[midway,below right=2pt,font=\small] {$\vec{u}=(u_1,u_2)$};
  \draw[quotientorange,very thick,->] (O) -- (V) node[midway,left=3pt,font=\small] {$\vec{v}=(v_1,v_2)$};

  \draw[structurepurple,thick,->] (O) -- (-0.33, 1.77) node[above left,font=\small] {$\vec{u}^\perp = (-u_2, u_1)$};

  \draw[->,gray!80] (0.7,0.13) arc (10.6:65.5:0.7);
  \node[font=\footnotesize,gray!80!black] at (0.85,0.55) {$\theta$};

  \node[targetgreen!75!black,font=\small\bfseries] at (2.6,1.4) {Area $= \norm{\vec{u}}\,h = \vec{u}^\perp\cdot\vec{v}$};
\end{tikzpicture}
\caption{The base-and-height derivation of the $2\times 2$ determinant. The parallelogram
spanned by $\vec{u}=(u_1,u_2)$ and $\vec{v}=(v_1,v_2)$ has base length $\norm{\vec{u}}$
and perpendicular altitude $h = \norm{\vec{v}}\sin\theta$. Rotating $\vec{u}$ by $90^\circ$
counterclockwise gives $\vec{u}^\perp = (-u_2, u_1)$, which points directly along the
altitude direction. The dot product $\vec{u}^\perp\cdot\vec{v} = u_1 v_2 - u_2 v_1$
matches base times height, producing $ad-bc$ when $\vec{u}=(a,c)$ and $\vec{v}=(b,d)$.}
\label{ch21:fig-base-height}
\end{figure}

\begin{formalbox}
For a $2\times 2$ matrix,
\[
  \det\begin{bmatrix} a & b \\ c & d \end{bmatrix} = ad - bc.
\]
Its absolute value is the area of the parallelogram spanned by the columns; its
sign is $+$ if the columns keep the same orientation as $(\vec e_1, \vec e_2)$
(counterclockwise) and $-$ if they reverse it.
\end{formalbox}

\begin{example}[Three quick $2\times 2$ determinants]
\label{ch21:ex-three-2x2}
\emph{A pure stretch.}
$\mat{A} = \begin{bmatrix} 2 & 0 \\ 0 & 3 \end{bmatrix}$ gives
$\det = 2\cdot 3 - 0\cdot 0 = 6$. Areas multiply by $6$: a $2$-by-$3$
rectangle, exactly as you'd expect from doubling width and tripling height.

\emph{A shear.}
$\mat{S} = \begin{bmatrix} 1 & 5 \\ 0 & 1 \end{bmatrix}$ gives
$\det = 1\cdot 1 - 5\cdot 0 = 1$. The shear slides the top of the square sideways
by $5$ but does \emph{not} change its area -- base $1$, height $1$, still area $1$.
This is the ``deck of cards'' fact made numerical.

\emph{A collapse.}
$\mat{C} = \begin{bmatrix} 1 & 2 \\ 2 & 4 \end{bmatrix}$ gives
$\det = 1\cdot 4 - 2\cdot 2 = 0$. The second column is twice the first, so both
columns lie on one line; the ``parallelogram'' is squashed flat onto that line and
has zero area. We will make a great deal of this $\det = 0$ case shortly.
\end{example}

\begin{labbox}
Compute these three $2\times 2$ determinants by hand, then \emph{draw} the
parallelogram of columns for each and check that the area matches:
\[
  \begin{bmatrix} 3 & 1 \\ 1 & 2 \end{bmatrix},\qquad
  \begin{bmatrix} 0 & -1 \\ 1 & 0 \end{bmatrix},\qquad
  \begin{bmatrix} 2 & 4 \\ 1 & 2 \end{bmatrix}.
\]
(Answers: $5$, $1$, $0$. The middle one is a $90^\circ$ rotation -- area
preserved, and we'll explain its positive sign in the next section. The last one
is singular.)
\end{labbox}

\section{Sign and Orientation: Tracking Reflections}
\label{ch21:sign}

So far we have discussed area, which is strictly non-negative. Where does the
\emph{sign} come from? It comes from orientation -- indicating whether the
transformation reflects space.

In the standard orientation of the plane, moving from the first basis vector
$\vec e_1$ to the second $\vec e_2$ proceeds \textbf{counterclockwise}. Now act
with $\mat{A}$ and inspect the image vectors $\mat{A}\vec e_1$ and $\mat{A}\vec
e_2$. If moving from the first image column to the second is still
counterclockwise, the transformation preserves orientation and $\det\mat{A} > 0$.
If the direction is now \emph{clockwise}, the transformation reflects the plane,
and $\det\mat{A} < 0$.

\begin{figure}[H]
\centering
\begin{tikzpicture}[scale=1.0,>=Stealth]
  % preserved
  \begin{scope}[xshift=0cm]
    \draw[->,gray!60] (-1.4,0) -- (1.6,0);
    \draw[->,gray!60] (0,-1.4) -- (0,1.6);
    \draw[vecblue,very thick,->] (0,0) -- (1.2,0.3) node[right]{$\mat{A}\vec e_1$};
    \draw[vecred,very thick,->] (0,0) -- (-0.4,1.1) node[above]{$\mat{A}\vec e_2$};
    \draw[->,targetgreen,thick] (0.7,0.18) arc (14:110:0.72);
    \node[targetgreen] at (0.55,0.95) {\small ccw};
    \node[below] at (0,-1.7) {$\det > 0$: orientation kept};
  \end{scope}
  % reversed
  \begin{scope}[xshift=6cm]
    \draw[->,gray!60] (-1.4,0) -- (1.6,0);
    \draw[->,gray!60] (0,-1.4) -- (0,1.6);
    \draw[vecblue,very thick,->] (0,0) -- (-0.4,1.1) node[above]{$\mat{A}\vec e_1$};
    \draw[vecred,very thick,->] (0,0) -- (1.2,0.3) node[right]{$\mat{A}\vec e_2$};
    \draw[->,trapred,thick] (-0.25,0.7) arc (110:14:0.72);
    \node[trapred] at (0.55,0.95) {\small cw};
    \node[below] at (0,-1.7) {$\det < 0$: orientation flipped};
  \end{scope}
\end{tikzpicture}
\caption{The sign of the determinant records orientation: does the transformation
keep the basis ordering counterclockwise (left, $\det>0$) or reverse it to
clockwise (right, $\det<0$)? A negative determinant indicates that the
transformation includes a reflection.}
\label{ch21:fig-orientation}
\end{figure}

\begin{intuitionbox}
A positive determinant corresponds to an orientation-preserving transformation:
stretches, rotations, and shears. A negative determinant indicates that the
transformation involves a \textbf{reflection}: the handedness of the coordinate
system is reversed. The magnitude $\abs{\det}$ tells you the factor by which
area changes, while the \emph{sign} indicates whether orientation is preserved
($+$) or reversed ($-$).
\end{intuitionbox}

\begin{example}[A reflection has determinant $-1$]
\label{ch21:ex-reflection}
The matrix that flips the plane across the $x$-axis is
$\mat{F} = \begin{bmatrix} 1 & 0 \\ 0 & -1 \end{bmatrix}$. It leaves $\vec e_1$
alone and sends $\vec e_2$ to $-\vec e_2$. Areas are clearly unchanged (a flip is
rigid), so $\abs{\det} = 1$; but orientation is reversed, so the sign is negative:
\[
  \det\mat{F} = (1)(-1) - (0)(0) = -1. \checkmark
\]
Compare the rotation by $90^\circ$, $\mat{R} = \begin{bmatrix} 0 & -1 \\ 1 & 0 \end{bmatrix}$,
which has $\det = (0)(0) - (-1)(1) = +1$: a rotation is rigid \emph{and} keeps
orientation, so its determinant is exactly $+1$. Reflection $-1$, rotation $+1$:
the sign is the whole difference between them.
\end{example}

\begin{trapbox}
A common slip is to think $\det\mat{A}$ is the area of the image and therefore
``always positive.'' It is the \emph{signed} area. The determinant $ad-bc$ can be
negative, and its sign records whether the transformation reverses orientation. Only $\abs{\det\mat{A}}$ is the (unsigned) area.
\end{trapbox}

\begin{recoverybox}
\textbf{Pause and think:} A reflection across a line in $2$D or across a plane in $3$D has determinant $-1$. What about a central inversion through the origin in $3$D, sending every vector $\vec{x}$ to $-\vec{x}$? Does it reverse orientation?
\emph{Answer:} In $3$D, reflecting through the origin is the matrix $\mat{A} = -\mat{I}_3 = \diag(-1,-1,-1)$. Because it reverses all three coordinate axes, each axis contributes a factor of $-1$, giving $\det(-\mat{I}_3) = (-1)^3 = -1$. It flips orientation: your right hand becomes a left hand, and no rigid physical rotation in $3$-space can undo that flip. (In $2$D, by contrast, $\det(-\mat{I}_2) = (-1)^2 = +1$: flipping both planar axes is simply a $180^\circ$ rotation, which preserves orientation.)
\end{recoverybox}

\begin{connectionbox}
Orientation connects linear algebra directly to multivariable calculus. When you change
variables in a multivariable integral, $\iint f\,dx\,dy$ becomes
$\iint f\,\abs{\det J}\,du\,dv$, where $J$ is the Jacobian matrix of the change of
coordinates. That factor $\abs{\det J}$ is the local area-scaling factor of
this chapter -- the change of variables formula applies the determinant-as-volume
concept to infinitesimal regions. The absolute value appears because a standard
scalar integral over a region depends only on the magnitude of the area element,
not on the orientation of the parameterization. In oriented integration
(differential forms and Stokes' theorem), the sign determines the orientation
of the boundary and domain.
\end{connectionbox}

\section{$\det = 0$ Means Collapse}
\label{ch21:singular}

Of all the values a determinant can take, one is special: zero. A zero determinant
is the geometric signature of \textbf{collapse}, and it is the single most useful
thing the determinant tells you.

If $\det\mat{A} = 0$, the unit square is sent to a parallelogram of zero area. But
a genuine parallelogram has positive area; the only way to get area zero is for the
two column vectors to lie on the \emph{same line} (in $2$D) or in the \emph{same
plane} (in $3$D). The machine has squashed a full-dimensional box down into
something thinner -- a line segment, or a flat sheet. Space lost a dimension.

\begin{figure}[H]
\centering
\begin{tikzpicture}[scale=1.0,>=Stealth]
  \begin{scope}
    \draw[->,gray!60] (-0.4,0) -- (2.6,0) node[right]{$x$};
    \draw[->,gray!60] (0,-0.4) -- (0,2.6) node[above]{$y$};
    \fill[vecblue!18] (0,0) rectangle (1,1);
    \draw[vecblue,very thick,->] (0,0) -- (1,0);
    \draw[vecblue,very thick,->] (0,0) -- (0,1);
    \node[below] at (0.5,-0.7) {area $1$};
  \end{scope}
  \node at (3.5,1.0) {\large $\xrightarrow{\;\mat{C}\;}$};
  \begin{scope}[xshift=5cm]
    \draw[->,gray!60] (-0.4,0) -- (3.2,0) node[right]{$x$};
    \draw[->,gray!60] (0,-0.4) -- (0,2.6) node[above]{$y$};
    % both columns on the line y=2x : (1,2) and (2,4)
    \draw[trapred,very thick,->] (0,0) -- (1,2) node[above left]{col $1$};
    \draw[trapred,very thick,->] (0,0) -- (1.4,2.8) node[right]{col $2$};
    \draw[trapred!50,dashed] (-0.1,-0.2) -- (1.5,3.0);
    \node[below] at (1.3,-0.7) {area $0$: a line};
  \end{scope}
\end{tikzpicture}
\caption{$\det\mat{C}=0$ for $\mat{C}=\protect\begin{bmatrix}1&2\\2&4\protect\end{bmatrix}$.
Both columns land on the single line $y=2x$, so the unit square is crushed flat to
a line segment of zero area. The machine is \emph{singular}: it has destroyed a
dimension and cannot be undone.}
\label{ch21:fig-collapse}
\end{figure}

This is the bridge between the geometry of this chapter and the algebra of
everything that came before. Collapse means the columns are \textbf{linearly
dependent}; equivalently, the matrix has rank less than full; equivalently, the
columns fail to span all of space; equivalently, $\mat{A}\vec x = \vec 0$ has a
nonzero solution (some direction gets crushed to the origin); equivalently,
$\mat{A}$ has no inverse. The determinant packages all of those conditions into a
single number you can check.

\begin{formalbox}
For a square matrix $\mat{A}$, the following are equivalent:
\begin{itemize}
  \item $\det\mat{A} = 0$;
  \item $\mat{A}$ is \textbf{singular} (not invertible);
  \item the columns of $\mat{A}$ are linearly dependent;
  \item $\rank\mat{A} < n$;
  \item the equation $\mat{A}\vec x = \vec 0$ has a nonzero solution;
  \item $\mat{A}$ collapses $\R^n$ onto a subspace of dimension $< n$.
\end{itemize}
Equivalently, $\det\mat{A}\neq 0$ iff $\mat{A}$ is invertible.
\end{formalbox}

\begin{aestheticbox}
This property is the primary reason determinants were studied centuries before
matrices were formalized. A single scalar, computed directly from the entries,
determines whether an entire linear system has a unique solution. The condition
$\det\neq 0$ unites several fundamental properties: the columns are linearly
independent, the transformation is invertible, and the null space contains only
the zero vector. Compressing a solvability test for $n$ equations into the
sign-or-zero of a single value rests on a geometric criterion: \emph{does the
transformation preserve volume, or does it collapse space into fewer
dimensions?}
\end{aestheticbox}

\begin{walkbox}
Picture a $3\times 3$ matrix whose three columns lie in a common plane through
the origin. The unit cube flattens into that plane, producing zero volume.
Mentally verify three consequences: the determinant is $0$; some nonzero
vector maps to $\vec 0$; and the original cube cannot be reconstructed from its
image, because distinct points in $\R^3$ collapse onto the same point in the
plane. Geometrically, this loss of dimension is why invertibility fails.
\end{walkbox}

\section{The $3\times 3$ Determinant: Volume of the Box}
\label{ch21:3x3}

In three dimensions the story is identical with ``area'' upgraded to
``volume.'' The unit cube has the three basis vectors $\vec e_1, \vec e_2, \vec
e_3$ as edges and volume $1$. A $3\times 3$ matrix sends it to the
\textbf{parallelepiped} (a slanted box) spanned by the three columns, and
$\abs{\det\mat{A}}$ is that box's volume. The sign now records ``right-handed''
versus ``left-handed'' -- whether the three image vectors obey the right-hand rule
or its mirror.

\begin{figure}[H]
\centering
\begin{tikzpicture}[scale=1.5,>=Stealth,line join=round]
  % a parallelepiped from three column vectors
  \coordinate (O) at (0,0,0);
  \coordinate (A) at (1.4,0.1,0.2);   % col1
  \coordinate (B) at (0.3,1.2,0.1);   % col2
  \coordinate (C) at (0.2,0.3,1.3);   % col3
  \coordinate (AB) at ($(A)+(B)$);
  \coordinate (AC) at ($(A)+(C)$);
  \coordinate (BC) at ($(B)+(C)$);
  \coordinate (ABC) at ($(A)+(B)+(C)$);
  % faces
  \fill[targetgreen!12] (O)--(A)--(AB)--(B)--cycle;
  \fill[targetgreen!20] (O)--(A)--(AC)--(C)--cycle;
  \fill[targetgreen!16] (O)--(B)--(BC)--(C)--cycle;
  \fill[targetgreen!10] (A)--(AB)--(ABC)--(AC)--cycle;
  \fill[targetgreen!10] (B)--(AB)--(ABC)--(BC)--cycle;
  \fill[targetgreen!8] (C)--(AC)--(ABC)--(BC)--cycle;
  % edges
  \draw[targetgreen,very thick,->] (O)--(A) node[below right]{col $1$};
  \draw[targetgreen,very thick,->] (O)--(B) node[left=2pt]{col $2$};
  \draw[targetgreen,very thick,->] (O)--(C) node[left]{col $3$};
  \draw[targetgreen!70] (A)--(AB) (B)--(AB) (A)--(AC) (C)--(AC) (B)--(BC) (C)--(BC);
  \draw[targetgreen!70] (AB)--(ABC) (AC)--(ABC) (BC)--(ABC);
  \node[above right=-2pt] at (AB) {volume $=\abs{\det\mat{A}}$};
\end{tikzpicture}
\caption{In $3$D the unit cube becomes the parallelepiped spanned by the three
columns of $\mat{A}$. Its volume is $\abs{\det\mat{A}}$, and the sign tells you
whether the three column vectors form a right-handed or left-handed frame.}
\label{ch21:fig-parallelepiped}
\end{figure}

To compute a $3\times 3$ determinant by hand, expand along the first row using
$2\times 2$ minors. For
\[
  \mat{A} = \begin{bmatrix} a & b & c \\ d & e & f \\ g & h & i \end{bmatrix},
\]
the rule (\emph{cofactor expansion}, which we'll justify thoroughly in the next
chapter) is
\[
  \det\mat{A}
  = a\det\!\begin{bmatrix} e & f \\ h & i \end{bmatrix}
  - b\det\!\begin{bmatrix} d & f \\ g & i \end{bmatrix}
  + c\det\!\begin{bmatrix} d & e \\ g & h \end{bmatrix}.
\]
Notice the alternating signs $+\,-\,+$: you cross out the row and column of each
top entry, take the determinant of what's left, and attach the checkerboard sign.
Written out,
\[
  \det\mat{A} = a(ei - fh) - b(di - fg) + c(dh - eg).
\]

\begin{example}[A worked $3\times 3$]
\label{ch21:ex-3x3}
Let
\[
  \mat{A} = \begin{bmatrix} 2 & 0 & 1 \\ 1 & 3 & 2 \\ 0 & 1 & 1 \end{bmatrix}.
\]
Expanding along the top row:
\[
  \det\mat{A}
  = 2\det\!\begin{bmatrix} 3 & 2 \\ 1 & 1 \end{bmatrix}
  - 0\cdot(\cdots)
  + 1\det\!\begin{bmatrix} 1 & 3 \\ 0 & 1 \end{bmatrix}.
\]
The two small determinants are $3\cdot1 - 2\cdot1 = 1$ and
$1\cdot1 - 3\cdot0 = 1$, so
\[
  \det\mat{A} = 2(1) - 0 + 1(1) = 3.
\]
The machine inflates volume by a factor of $3$, and the positive sign says it keeps
the right-handed orientation of space. You can expand along any row or column you
like and get the same answer, so it pays to pick the one with the most zeros: here
the first column or third row, each with a zero, would be just as fast.
\end{example}

\begin{trapbox}
The ``diagonals'' trick (Sarrus' rule) -- multiply the three downhill diagonals,
subtract the three uphill diagonals -- works \textbf{only for $3\times 3$}. There
is no $4\times 4$ version of it; applying it in higher dimensions yields incorrect
results. For larger matrices, use cofactor expansion or row reduce to triangular
form (next section) and multiply the pivots. Sarrus' rule applies only to
$3\times 3$ matrices, not to general dimensions.
\end{trapbox}

\section{How Row Operations Change the Determinant}
\label{ch21:row-ops}

Computing large determinants via cofactor expansion is computationally prohibitive:
an $n\times n$ expansion involves $n!$ terms, which for $n = 20$ exceeds
$2.4\times 10^{18}$ arithmetic operations. The practical method uses
\textbf{Gaussian elimination}. We simply track how each elementary row operation
affects the determinant. There are three operations, each with a clear
geometric effect.

\textbf{(1) Add a multiple of one row to another: determinant unchanged.}
This is the shear. Adding row $i$ to row $j$ slides one edge of the box parallel to
another; base and height are unchanged, so the volume -- and hence the
determinant -- is unchanged. This is the primary tool: you may do as many of these as
you like for free.

\textbf{(2) Swap two rows: determinant flips sign.}
Swapping two rows swaps two edges of the box, which reverses orientation -- the
right hand becomes a left hand. The volume is the same but the sign flips:
$\det \to -\det$.

\textbf{(3) Multiply a row by a scalar $k$: determinant multiplies by $k$.}
Stretching one edge of the box by $k$ stretches the volume by $k$. So scaling a
single row by $k$ scales the determinant by $k$.

\begin{figure}[H]
\centering
\begin{tikzpicture}[scale=0.9,>=Stealth]
  % shear: no change
  \begin{scope}[xshift=0cm]
    \fill[gray!15] (0,0) rectangle (1.2,1.2);
    \fill[targetgreen!20] (0,0)--(1.2,0)--(1.8,1.2)--(0.6,1.2)--cycle;
    \draw[targetgreen,thick] (0,0)--(1.2,0)--(1.8,1.2)--(0.6,1.2)--cycle;
    \node[below] at (0.9,-0.5) {(1) shear};
    \node at (0.9,-1.25) {\small $\det$ unchanged};
  \end{scope}
  % swap: sign flip
  \begin{scope}[xshift=4.2cm]
    \draw[->,vecblue,thick] (0,0)--(1.2,0) node[right]{\small$1$};
    \draw[->,vecred,thick] (0,0)--(0,1.2) node[above]{\small$2$};
    \draw[->,gray] (2.0,0.6)--(2.7,0.6);
    \draw[->,vecred,thick] (3.4,0)--(4.6,0) node[right]{\small$2$};
    \draw[->,vecblue,thick] (3.4,0)--(3.4,1.2) node[above]{\small$1$};
    \node[below] at (2.3,-0.5) {(2) swap};
    \node at (2.3,-1.25) {\small $\det \to -\det$};
  \end{scope}
  % scale
  \begin{scope}[xshift=11cm]
    \fill[targetgreen!20] (0,0) rectangle (1.2,2.0);
    \draw[targetgreen,thick] (0,0) rectangle (1.2,2.0);
    \draw[<->,gray] (1.5,0)--(1.5,2.0) node[midway,right]{\small$\times k$};
    \node[below] at (0.6,-0.5) {(3) scale row};
    \node at (0.6,-1.25) {\small $\det \to k\det$};
  \end{scope}
\end{tikzpicture}
\caption{The three elementary row operations and their effect on the determinant.
Shears (add a multiple of a row) leave it alone, swaps flip its sign, and scaling a
row by $k$ multiplies it by $k$.}
\label{ch21:fig-rowops}
\end{figure}

These three rules turn determinant computation into elimination. Row reduce
$\mat{A}$ to an upper-triangular matrix $\mat{U}$, keeping a running tally: count
the swaps (each contributes a factor $-1$) and record any row-scalings. Then use
one more beautiful fact.

\begin{theorem}[Determinant of a triangular matrix]
\label{ch21:thm-triangular}
If $\mat{T}$ is (upper or lower) triangular, then $\det\mat{T}$ is the product of
its diagonal entries.
\end{theorem}

\begin{proof}[Why it's true, geometrically and by expansion]
Geometrically: a triangular matrix's columns are already ``staircased,'' and the
box they span is a sheared rectangular box whose edges along each axis have lengths
equal to the diagonal entries; shears don't change volume, so the volume is just
the product of those edge lengths. Algebraically: take an upper-triangular matrix
and expand along its first \emph{column}. In an upper-triangular matrix the first
column has only the top-left entry $t_{11}$ nonzero (everything below the diagonal
is zero), so only the $(1,1)$ cofactor survives, leaving $t_{11}$ times the
determinant of the smaller upper-triangular matrix obtained by deleting row $1$ and
column $1$. Induct on that smaller matrix, and you peel off
$t_{11}t_{22}\cdots t_{nn}$. (For a lower-triangular matrix, expand along the first
\emph{row} instead; the argument is the mirror image.)
\end{proof}

\begin{example}[Determinant by elimination]
\label{ch21:ex-elim}
Compute $\det\mat{A}$ for
\[
  \mat{A} = \begin{bmatrix} 2 & 1 & 1 \\ 4 & 3 & 1 \\ -2 & 1 & 5 \end{bmatrix}.
\]
Eliminate below the first pivot with shears (free, $\det$ unchanged):
$R_2 \leftarrow R_2 - 2R_1$ and $R_3 \leftarrow R_3 + R_1$:
\[
  \begin{bmatrix} 2 & 1 & 1 \\ 0 & 1 & -1 \\ 0 & 2 & 6 \end{bmatrix}.
\]
Now $R_3 \leftarrow R_3 - 2R_2$ (again free):
\[
  \mat{U} = \begin{bmatrix} 2 & 1 & 1 \\ 0 & 1 & -1 \\ 0 & 0 & 8 \end{bmatrix}.
\]
No swaps, no scalings, so $\det\mat{A} = \det\mat{U} = 2\cdot 1\cdot 8 = 16$. The
diagonal entries of $\mat{U}$ are exactly the pivots of elimination -- which is why
people say \emph{the determinant is the product of the pivots} (times $(-1)^{\#
\text{swaps}}$). If at any point a pivot had been forced to $0$, the determinant
would be $0$ and the matrix singular.
\end{example}

\begin{strangbox}
Strang builds the determinant axiomatically: he lists three defining properties --
$\det I = 1$, the sign flips under a row swap, and the determinant is linear in
each row separately -- and shows that those three force the formula $ad-bc$ and
everything else. That is the elegant top-down route, and after Strang you should
absolutely go back and admire how those three axioms generate the whole theory
(rules (1), (2), (3) above are corollaries of them). We have taken a geometric
approach instead: volume first, formulas second, because volume scaling provides
the intuition behind Strang's axioms. ``Linear in each row'' corresponds to
stretching one edge of the box; ``swap flips the sign'' reflects reversal of
orientation. Both viewpoints develop the same underlying theory.
\end{strangbox}

\begin{labbox}
Compute $\det\begin{bmatrix} 1 & 2 & 3 \\ 2 & 5 & 3 \\ 1 & 0 & 8 \end{bmatrix}$
twice: once by cofactor expansion along the first row, and once by elimination to
triangular form. You should get $-1$ both times. If your two answers disagree, the
likeliest culprit is a forgotten sign in the cofactor checkerboard or an
uncounted row swap.
\end{labbox}

\section{The Product Rule: \texorpdfstring{$\det(\mat{A}\mat{B}) = \det\mat{A}\,\det\mat{B}$}{det(AB)=detA detB}}
\label{ch21:product}

The product rule follows directly from the volume interpretation.
Recall that the product $\mat{A}\mat{B}$ corresponds to the composition ``apply $\mat{B}$
first, then $\mat{A}$.'' Apply that composite mapping to the unit box:

\begin{itemize}
  \item First $\mat{B}$ scales the box's volume by $\det\mat{B}$.
  \item Then $\mat{A}$ scales \emph{that} volume by $\det\mat{A}$.
\end{itemize}

Scaling factors of successive operations multiply (Figure~\ref{ch21:fig-product}),
so the total scaling is $\det\mat{A}\cdot\det\mat{B}$. And the total scaling of the
combined machine $\mat{A}\mat{B}$ is by definition $\det(\mat{A}\mat{B})$.
Therefore:

\begin{figure}[H]
\centering
\begin{tikzpicture}[scale=1.0,>=Stealth]
  % start: unit box
  \begin{scope}[xshift=0cm]
    \fill[vecblue!18] (0,0) rectangle (0.8,0.8);
    \draw[vecblue,thick] (0,0) rectangle (0.8,0.8);
    \node[below] at (0.4,-0.4) {\small box, vol $1$};
  \end{scope}
  \node at (1.9,0.4) {\large $\xrightarrow{\;\mat{B}\;}$};
  % after B: scaled by det B
  \begin{scope}[xshift=3.0cm]
    \fill[quotientorange!22] (0,0)--(1.3,0)--(1.7,0.9)--(0.4,0.9)--cycle;
    \draw[quotientorange,thick] (0,0)--(1.3,0)--(1.7,0.9)--(0.4,0.9)--cycle;
    \node[below] at (0.7,-0.4) {\small vol $\det\mat{B}$};
  \end{scope}
  \node at (5.5,0.4) {\large $\xrightarrow{\;\mat{A}\;}$};
  % after A: scaled again by det A
  \begin{scope}[xshift=6.7cm]
    \fill[targetgreen!22] (0,0)--(2.0,0.3)--(2.5,1.5)--(0.5,1.2)--cycle;
    \draw[targetgreen,thick] (0,0)--(2.0,0.3)--(2.5,1.5)--(0.5,1.2)--cycle;
    \node[below] at (1.1,-0.4) {\small vol $\det\mat{A}\cdot\det\mat{B}$};
  \end{scope}
\end{tikzpicture}
\caption{The product rule, seen as composition. Doing $\mat{B}$ first multiplies
the box's volume by $\det\mat{B}$; then $\mat{A}$ multiplies that volume by
$\det\mat{A}$. The combined machine $\mat{A}\mat{B}$ therefore multiplies volume by
$\det\mat{A}\,\det\mat{B}$, which is exactly $\det(\mat{A}\mat{B})$.}
\label{ch21:fig-product}
\end{figure}

\begin{formalbox}
For any two $n\times n$ matrices,
\[
  \det(\mat{A}\mat{B}) = \det\mat{A}\,\det\mat{B}.
\]
The determinant turns ``compose machines'' (matrix multiplication) into ``multiply
numbers.'' Orientation behaves correctly too: two flips ($\det\mat{A}<0$,
$\det\mat{B}<0$) compose to no flip ($\det(\mat{A}\mat{B})>0$), exactly as
$(-)(-)=+$.
\end{formalbox}

\begin{intuitionbox}
Volume scalings \emph{compose by multiplication}. If $\mat{B}$ triples volume and
$\mat{A}$ doubles it, then doing $\mat{B}$ then $\mat{A}$ multiplies volume by
$6$. That is the entire content of $\det(\mat{A}\mat{B}) = \det\mat{A}\,\det\mat{B}$
-- a statement that would be a nightmare to prove by grinding through the entry
formulas, but is a one-line truth about scaling factors.
\end{intuitionbox}

\begin{example}[Checking the product rule by hand]
\label{ch21:ex-product}
Let
\[
  \mat{A} = \begin{bmatrix} 2 & 1 \\ 0 & 3 \end{bmatrix},\qquad
  \mat{B} = \begin{bmatrix} 1 & 4 \\ 2 & 1 \end{bmatrix}.
\]
Then $\det\mat{A} = 6$ and $\det\mat{B} = 1 - 8 = -7$, so the product rule
predicts $\det(\mat{A}\mat{B}) = 6\cdot(-7) = -42$. Check directly:
\[
  \mat{A}\mat{B}
  = \begin{bmatrix} 2 & 1 \\ 0 & 3 \end{bmatrix}\begin{bmatrix} 1 & 4 \\ 2 & 1 \end{bmatrix}
  = \begin{bmatrix} 4 & 9 \\ 6 & 3 \end{bmatrix},
  \qquad
  \det(\mat{A}\mat{B}) = 12 - 54 = -42.\ \checkmark
\]
\end{example}

The product rule has consequences that fall out instantly -- each a small
geometric story.

\begin{corollary}[Determinant of an inverse]
\label{ch21:cor-inverse}
If $\mat{A}$ is invertible, then $\det(\mat{A}\inv) = \dfrac{1}{\det\mat{A}}$.
\end{corollary}

\begin{proof}
Apply the product rule to $\mat{A}\mat{A}\inv = \mat{I}$. Since $\det\mat{I} = 1$
(the identity changes no volume),
\[
  \det\mat{A}\cdot\det(\mat{A}\inv) = \det\mat{I} = 1,
\]
so $\det(\mat{A}\inv) = 1/\det\mat{A}$. Geometrically: if $\mat{A}$ multiplies
volume by $5$, its inverse must divide volume by $5$ to put everything back. And
notice this proof \emph{requires} $\det\mat{A}\neq 0$ -- you cannot divide by
zero, which is the algebraic version of ``a collapsed box can't be un-collapsed.''
\end{proof}

\begin{corollary}[Powers and similar matrices]
\label{ch21:cor-powers}
$\det(\mat{A}^k) = (\det\mat{A})^k$ for any integer $k\ge 1$, and similar matrices
have equal determinants: $\det(\mat{P}\inv\mat{A}\mat{P}) = \det\mat{A}$.
\end{corollary}

\begin{proof}
The first is the product rule applied $k$ times. For the second,
\[
  \det(\mat{P}\inv\mat{A}\mat{P})
  = \det(\mat{P}\inv)\,\det\mat{A}\,\det\mat{P}
  = \frac{1}{\det\mat{P}}\,\det\mat{A}\,\det\mat{P}
  = \det\mat{A}.
\]
This is why the
determinant is an \emph{intrinsic} property of a linear map, independent of which
basis you write it in: changing basis is exactly conjugation by $\mat{P}$, and the
volume-scaling of a map doesn't care what coordinates you use to describe it.
\end{proof}

\begin{remark}[A warning about sums]
\label{ch21:remark-sum}
The determinant multiplies beautifully but does \emph{not} add. In general
$\det(\mat{A}+\mat{B}) \neq \det\mat{A} + \det\mat{B}$. For instance with
$\mat{A} = \mat{B} = \mat{I}$ in $\R^2$, $\det(\mat{A}+\mat{B}) = \det(2\mat{I})
= 4$, while $\det\mat{A} + \det\mat{B} = 1 + 1 = 2$. There is no ``volume of a
sum'' -- geometrically, adding matrices does not correspond to any clean operation
on the boxes, whereas multiplying them does (compose the machines).
\end{remark}

\begin{connectionbox}
\textbf{Gram--Schmidt and QR: volume as a product of altitudes.}
In Chapter~20, we decomposed any matrix $\mat{A}$ with independent columns into
$\mat{A} = \mat{Q}\mat{R}$, where $\mat{Q}$ has orthonormal columns and $\mat{R}$ is
upper triangular. Now we can see what that factorization reveals about volume. By the
product rule,
\[
  \abs{\det\mat{A}} = \abs{\det\mat{Q}}\,\abs{\det\mat{R}}.
\]
Because $\mat{Q}$ is an orthogonal matrix ($\mat{Q}\T\mat{Q} = \mat{I}$), it preserves
lengths and angles -- a rigid rotation or reflection -- so $\abs{\det\mat{Q}} = 1$. And
because $\mat{R}$ is upper triangular, its determinant is simply the product of its
diagonal entries: $\det\mat{R} = r_{11} r_{22} \cdots r_{nn}$. Therefore,
\[
  \abs{\det\mat{A}} = r_{11} r_{22} \cdots r_{nn}.
\]
Look at what this means geometrically: $r_{11} = \norm{\vec{a}_1}$ is the length of the
first edge (the base). The entry $r_{22}$ is the perpendicular altitude of the second
column above the first. The entry $r_{33}$ is the perpendicular altitude of the third
column above the plane of the first two, and so on. The determinant is literally
$\text{base} \times \text{height}_1 \times \text{height}_2 \times \cdots \times \text{height}_{n-1}$.
Gram--Schmidt's orthogonalization was building the determinant all along.
\end{connectionbox}

\begin{recoverybox}
\textbf{Pause and think:} If $\mat{A}$ is a $3\times 3$ matrix and you double every entry to form $2\mat{A}$, does its determinant double?
\emph{Answer:} No, it multiplies by $8 = 2^3$. Doubling the matrix scales all three edges of the parallelepiped by $2$, so the volume scales by $2\times 2\times 2 = 8$. In general, because $k\mat{A} = (k\mat{I})\mat{A}$, the product rule gives $\det(k\mat{A}) = \det(k\mat{I})\det\mat{A} = k^n\det\mat{A}$ for an $n\times n$ matrix.
\end{recoverybox}

\begin{applicationbox}
\textbf{Why generative models care about determinants.} A \emph{normalizing flow}
is a modern machine-learning model that learns a probability distribution by
warping a simple one (a Gaussian blob) through an invertible neural network
$f$. The change-of-variables formula says that if a point $\vec z$ has density
$p_Z(\vec z)$, then its image $\vec x = f(\vec z)$ has density
\[
  p_X(\vec x) = p_Z(\vec z)\,\bigl|\det J_f(\vec z)\bigr|^{-1},
\]
where $J_f$ is the Jacobian of $f$. The $\abs{\det J_f}$ is doing \emph{exactly}
the job from this chapter: it is the local volume-scaling factor of the warp, and
it has to appear because probability mass is conserved -- if the map stretches a
region to twice its volume, the density there must drop by half so the total still
integrates to $1$. Training such a model means computing (and differentiating)
$\log\abs{\det J_f}$ millions of times, so flow architectures are designed with
triangular Jacobians on purpose -- because, as you now know, a triangular
determinant is just the product of the diagonal, cheap to compute. The entire
field of normalizing flows is an engineering response to ``make the determinant
fast.''
\end{applicationbox}

\begin{historybox}
Today we introduce determinants as functions of square matrices, but historically the sequence was reversed: determinants were discovered more than a century and a half \emph{before} anyone defined a matrix. In the late seventeenth century, mathematicians faced an immediate practical roadblock: how to tell whether a large system of linear equations has a unique solution without slogging through pages of variable elimination. In Japan, Seki Takakazu (1683) worked out systematic tables of elimination coefficients up to five variables using bamboo counting rods (\emph{sangi}, traditional sticks arranged on a counting board to represent numerical coefficients). A decade later in Europe, Gottfried Wilhelm Leibniz (1693) independently derived equivalent solvability conditions for three equations in two unknowns in a letter to Guillaume de l'H\^opital.

For over a century, however, determinants remained confined to coordinate algebra, formulated as polynomials of $n!$ signed terms (all $n!$ ways to pair rows with columns, which for a modest $10\times 10$ system already exceeds $3.6$ million terms). In that coordinate language, proving basic facts like the product rule $\det(\mat{A}\mat{B}) = \det\mat{A}\,\det\mat{B}$ required tedious algebraic substitutions. The breakthrough came when Joseph-Louis Lagrange (1773) recognized that a $3\times 3$ determinant calculates the volume of a tetrahedron (a four-sided triangular pyramid), and Augustin-Louis Cauchy (1812) unified the subject by proving that the determinant is the signed volume scaling factor of space, establishing the modern name and notation. Carl Gustav Jacob Jacobi (1841) later extended this to calculus via the Jacobian (the matrix of local partial derivatives measuring how coordinate transformations stretch infinitesimal volume elements and flip orientation).

This geometric shift created a nineteenth-century debate. When Arthur Cayley and James Joseph Sylvester introduced the concept of a \emph{matrix} in the 1850s, many traditional algebraists resisted, arguing that determinants were the primary objects of study and matrices were merely auxiliary notation. The dispute was settled only when linear transformations took center stage: matrices describe the linear mapping, while the determinant provides one scalar summary of that mapping.

What this theory provided was a single scalar test for invertibility ($\det\mat{A} \neq 0$) and a precise tracker of whether a linear transformation preserves spatial handedness or reflects it. Collapsing an $n\times n$ transformation to a single scalar, however, discards directional information. A matrix can stretch one axis by a factor of a million and compress another by a factor of a millionth while leaving the volume scaling exactly $1$. In addition, calculating determinants via $n!$ permutation terms is computationally infeasible for large matrices, and floating-point roundoff can underflow a small product to zero. In Chapter~22 we develop cofactors and Cramer's rule to study symbolic inverses, and in Chapter~29 we meet the singular value decomposition, which restores the individual stretch factors that the determinant aggregates together.
\end{historybox}

\section{Putting It Together}
\label{ch21:together}

The geometry of the transformed unit box explains the key properties of the determinant:

\begin{itemize}
  \item The determinant is the \textbf{signed volume-scaling factor} of a linear
        map. Its magnitude is how much volume changes; its sign is whether
        orientation is preserved ($+$) or reversed ($-$).
  \item In $2$D, $\det\begin{bmatrix} a & b \\ c & d \end{bmatrix} = ad - bc$, the
        signed area of the parallelogram of columns. In $3$D, cofactor expansion
        gives the signed volume of the parallelepiped of columns.
  \item $\det\mat{A} = 0$ exactly when the box collapses -- when columns are
        dependent, the map is singular, and no inverse exists. This is the
        determinant's most important job.
  \item Row operations have clean effects: shears leave $\det$ alone, swaps flip
        its sign, scalings multiply it. So $\det = (-1)^{\#\text{swaps}}$ times the
        product of the pivots -- the practical way to compute.
  \item $\det(\mat{A}\mat{B}) = \det\mat{A}\,\det\mat{B}$, because volume scalings
        compose by multiplication. From this fall $\det(\mat{A}\inv) =
        1/\det\mat{A}$ and the basis-independence of the determinant.
\end{itemize}

\begin{aestheticbox}
The determinant is the first place in linear algebra where a single scalar
measures how a geometric region transforms under a linear mapping. This concept
reappears throughout subsequent chapters: as the invertibility criterion, in
Cramer's rule and cofactor formulas, in the characteristic polynomial
$\det(\mat{A} - \lambda\mat{I})$ used to find eigenvalues, and as the product of
singular values in the SVD. Each application traces back to the same geometric
question: \emph{by what factor does the transformation scale volume, and does it
preserve orientation?} Viewing the determinant through volume clarifies why
these algebraic formulas take the form they do.
\end{aestheticbox}

We now understand what the determinant means: it is signed volume scaling. But two
pressing practical questions remain. First, how do we peel apart a high-dimensional
determinant into smaller, manageable sub-determinants without running elimination?
Second, can we harness the determinant to write down an explicit, closed-form
formula for the inverse matrix $\mat{A}\inv$, and for the exact solution $\vec{x}$
of a linear system $\mat{A}\vec{x}=\vec{b}$? In Chapter~22, we answer both questions
by building cofactors, the adjugate matrix, and Cramer's rule.

\section{Exercises}
\label{ch21:exercises}

\begin{exercise}[Mental checks: no pencil needed]
\label{ch21:ex-mental}
Answer each of these in five seconds by visualizing the geometry:
\begin{enumerate}[label=(\alph*)]
  \item If $\mat{A}$ is a $4\times 4$ matrix with $\det\mat{A} = 3$, what is $\det(2\mat{A})$?
  \item A $3\times 3$ matrix has a row of zeros. What is its determinant, and what happened to the unit cube?
  \item If you swap two columns of an $n\times n$ matrix, what happens to the determinant?
  \item If $\det\mat{A} = 5$ and $\det\mat{B} = -2$, what is $\det(\mat{A}^2\mat{B}\inv)$?
  \item Can an orthogonal matrix have determinant equal to $3$? Why or why not?
\end{enumerate}
\end{exercise}

\begin{exercise}[Warm-up areas]
\label{ch21:ex-warmup}
Compute the determinant of each, and state the area (or volume) and the
orientation (preserved or reversed):
\[
  \text{(a)}\ \begin{bmatrix} 4 & 2 \\ 1 & 3 \end{bmatrix},\quad
  \text{(b)}\ \begin{bmatrix} 1 & 6 \\ 0 & 1 \end{bmatrix},\quad
  \text{(c)}\ \begin{bmatrix} 0 & 2 \\ 3 & 0 \end{bmatrix},\quad
  \text{(d)}\ \begin{bmatrix} 3 & 6 \\ 1 & 2 \end{bmatrix}.
\]
\end{exercise}

\begin{exercise}[Draw the parallelogram]
\label{ch21:ex-draw}
For $\mat{A} = \begin{bmatrix} 2 & 1 \\ 1 & 3 \end{bmatrix}$, draw the unit square
and its image under $\mat{A}$ on graph paper. Measure the area of the image
parallelogram by the surrounding-rectangle-minus-triangles method, and confirm it
equals $\det\mat{A}$.
\end{exercise}

\begin{exercise}[A $3\times 3$ two ways]
\label{ch21:ex-3x3-two-ways}
Compute $\det\begin{bmatrix} 1 & 2 & 0 \\ 3 & -1 & 2 \\ 2 & 0 & 1 \end{bmatrix}$
first by cofactor expansion along whichever row or column has the most zeros, then
again by elimination to triangular form. Confirm the answers match.
\end{exercise}

\begin{exercise}[When is it singular?]
\label{ch21:ex-singular}
Find all values of $t$ for which
$\mat{A} = \begin{bmatrix} t & 1 & 0 \\ 1 & t & 1 \\ 0 & 1 & t \end{bmatrix}$
is singular. (Set up $\det\mat{A} = 0$ and solve the resulting cubic in $t$.)
Interpret your answer geometrically: for those $t$, what has happened to the box?
\end{exercise}

\begin{exercise}[Row operations, tracked]
\label{ch21:ex-rowops}
Start from $\mat{A} = \begin{bmatrix} 0 & 2 & 1 \\ 1 & 1 & 0 \\ 3 & 0 & 2 \end{bmatrix}$.
Reduce it to upper-triangular form, carefully recording each swap and each scaling.
Use your record plus Theorem~\ref{ch21:thm-triangular} to compute $\det\mat{A}$, and
state how many sign flips you incurred.
\end{exercise}

\begin{exercise}[Product rule by hand]
\label{ch21:ex-product-hand}
Let $\mat{A} = \begin{bmatrix} 1 & 2 \\ 3 & 4 \end{bmatrix}$ and
$\mat{B} = \begin{bmatrix} 2 & 0 \\ 1 & 5 \end{bmatrix}$. Compute $\det\mat{A}$,
$\det\mat{B}$, the product $\mat{A}\mat{B}$, and $\det(\mat{A}\mat{B})$ directly.
Verify $\det(\mat{A}\mat{B}) = \det\mat{A}\,\det\mat{B}$. Then compute
$\det(\mat{A}\inv)$ without forming $\mat{A}\inv$.
\end{exercise}

\begin{exercise}[Scaling the whole matrix]
\label{ch21:ex-scaling}
For an $n\times n$ matrix $\mat{A}$ and a scalar $k$, prove that
$\det(k\mat{A}) = k^n\det\mat{A}$. (Hint: $k\mat{A}$ scales \emph{every} one of the
$n$ rows by $k$; apply the row-scaling rule $n$ times.) Then explain geometrically
why the exponent is $n$ and not $1$: what happens to the volume of an
$n$-dimensional box when you scale all $n$ edges by $k$?
\end{exercise}

\begin{exercise}[Orientation of a reflection]
\label{ch21:ex-reflection-line}
Let $\mat{F}_\theta$ reflect the plane across the line through the origin at angle
$\theta$. Its matrix is
$\mat{F}_\theta = \begin{bmatrix} \cos 2\theta & \sin 2\theta \\ \sin 2\theta & -\cos 2\theta \end{bmatrix}$.
Compute $\det\mat{F}_\theta$ and show it equals $-1$ for \emph{every} $\theta$.
Explain in one sentence why this had to be the answer before you computed
anything.
\end{exercise}

\begin{exercise}[A determinant that proves dependence]
\label{ch21:ex-dependence}
Three vectors $\vec a = \colvec{1\\2\\3}$, $\vec b = \colvec{0\\1\\4}$,
$\vec c = \colvec{2\\5\\10}$ in $\R^3$. By computing a single $3\times 3$
determinant, decide whether they are linearly independent. If they are dependent,
find scalars (not all zero) with $x\vec a + y\vec b + z\vec c = \vec 0$, and relate
your finding to the determinant being zero.
\end{exercise}

\begin{exercise}[Challenge: the area of a triangle]
\label{ch21:ex-triangle}
Show that the area of the triangle with vertices $\vec p, \vec q, \vec r$ in the
plane is $\tfrac{1}{2}\bigl|\det[\,\vec q - \vec p \mid \vec r - \vec p\,]\bigr|$,
where the $2\times 2$ matrix has columns $\vec q - \vec p$ and $\vec r - \vec p$.
Use it to find the area of the triangle with vertices $(1,1)$, $(4,2)$, $(2,5)$.
(Hint: the triangle is half the parallelogram spanned by the two edge vectors
from $\vec p$.)
\end{exercise}

\begin{exercise}[Challenge: a proof of the product rule for the $2\times 2$ case]
\label{ch21:ex-product-proof}
Without invoking the geometric argument, prove $\det(\mat{A}\mat{B}) =
\det\mat{A}\,\det\mat{B}$ for general $2\times 2$ matrices by brute force: write
$\mat{A} = \begin{bmatrix} a & b \\ c & d \end{bmatrix}$,
$\mat{B} = \begin{bmatrix} p & q \\ r & s \end{bmatrix}$, multiply them out, and
expand both $\det(\mat{A}\mat{B})$ and $(ad-bc)(ps-qr)$. Watch the cross terms
cancel. (This is good algebraic exercise -- and it should make you grateful for the
one-line geometric proof.)
\end{exercise}
